\documentclass{article}
\usepackage[T1]{fontenc}
\usepackage{lmodern}
\usepackage{graphicx}
\usepackage{amsmath,amssymb}
\usepackage[a4paper,margin=2cm]{geometry}
\usepackage[absolute,overlay]{textpos}
\setlength{\TPHorizModule}{1mm}
\setlength{\TPVertModule}{1mm}
\usepackage[indonesian]{babel}
\usepackage{hyperref}
\setlength{\emergencystretch}{2em}
% unit: Halaman 30

\begin{document}

% === Halaman 30 (llm) ===

\begin{itemize}
    \item Contoh: $\text{ROT13}(\text{ROTATE}) = \text{EBGNGR}$
    \item Nama ``ROT13'' berasal dari \textit{net.jokes} (\textit{http://groups.google.com/group/net.jokes}) (tahun 1980)
    \item ROT13 biasanya digunakan di dalam forum \textit{online} untuk menyandikan jawaban teka-teki, kuis, canda, dsb
    \item Enkripsi arsip dua kali dengan $\text{ROT13}$ menghasilkan pesan semula:
    \[ P = \text{ROT13}(\text{ROT13}(P)) \]
    sebab $\text{ROT}_{13}(\text{ROT}_{13}(x)) = \text{ROT}_{26}(x) = x$
    \item Jadi dekripsi cukup dilakukan dengan mengenkripsi cipherteks kembali dengan $\text{ROT13}$
\end{itemize}

\end{document}
